% Pure insertion fragment for the existing original report. No preamble.
% Theory and implemented operations only; development-result paragraph is
% appended after the complete independent review is available.
\subsection{EV 的可执行灵活性：曲线、地点和人员接入}
\label{evf:flexibility}

本轮把设备优势与算法优势分开检查。设备优势来自更多可执行的接入和功率选择；算法优势要求优化确实利用这些选择，并改善新增计量后的完整 AC 竞争拟合。设计仍使用式\eqref{si:robust-task}的局部信息增量，最终评价使用所有候选共同解释历史和新增数据的实际似然差。

\subsubsection{同一个充电服务，怎样公平比较 EV 与 ESS}

令 $u_t^+,u_t^-\ge0$ 分别为网侧充、放电功率，单位 kW，$u_t=u_t^+-u_t^-$；$e_t$ 为电池能量，单位 kWh，$\Delta t$ 为时长，单位小时。充、放电效率分别为 $\eta_c,\eta_d\in(0,1]$。移动到地点 $j$ 的驾驶耗电为 $E_{{\rm drive},j}$，出发前电量为 $e_{\rm initial}$。接入后的状态满足
\begin{equation}
 \boxed{e_0=e_{\rm initial}-E_{{\rm drive},j},\qquad
 e_{t+1}=e_t+\Delta t(\eta_cu_t^+-u_t^-/\eta_d).}
 \label{evf:battery}
\end{equation}
各时段还要满足功率上限、电量上下限、可接入窗口和充放电互斥。普通单向设施取 $u_t^-=0$。本轮固定相同的精确离站电量 $e_T=e_{\rm depart}$，另检查返程后电量和备用要求；本批返程耗电设为零。未来返程耗电因地点不同时，应按共同返程后服务电量调整任务，或明确报告该差异。

令 $C=\Delta t\sum_tu_t^+$、$D=\Delta t\sum_tu_t^-$ 为总充、放电量，单位 kWh。对式\eqref{evf:battery}求和得到
\begin{equation}
 \boxed{C-D=e_T-e_{\rm initial}+E_{{\rm drive},j}
 +(1-\eta_c)C+(1/\eta_d-1)D.}
 \label{evf:energy-account}
\end{equation}
等式右侧依次是电池服务增量、驾驶能量和转换损耗。因此固定末端电量后，V2G 可以比只充电消耗更多网侧净电量；吞吐量 $C+D$ 也须单列。若反而同时固定为单向任务的网侧电量，并要求相同末端电量，在 $\eta_c\eta_d<1$ 时上式迫使 $D=0$。本轮比较固定电池服务，另外记录网侧电量、吞吐、损耗和移动。

同地点、相同实际 P/Q 曲线、背景、AC 分支及相同计量和功率追踪误差模型下，EV 与 ESS 的 AC 辨识响应相同。证明：设备名称不进入 AC 方程；对每个候选，节点注入和其余物理输入相同即给出相同计量分布，因而残差、似然函数和基于分布的辨识指标相同。比较局部求解器已取得值时，还应使用相同拟合规则和随机状态。EV 的实际区别是可安排的用车时间、支持时的双向功率，以及能够到达的接入地点；固定 ESS 的优势是较稳定的接入条件。相同总功率放在不同节点时，保留空间注入 $p_{j,t}$，不能只合并成一条总曲线。

\subsubsection{设备自由度增加，何时保证有更多可用信息}

上述操作约束定义可行曲线集合 $\mathcal F_j$。若 $\mathcal F_1\subseteq\mathcal F_2$，两者使用相同的 $\Phi$，且最优值存在，则
\begin{equation}
 \max_{u\in\mathcal F_2}\Phi(u)\ge
 \max_{u\in\mathcal F_1}\Phi(u).
 \label{evf:inclusion}
\end{equation}
证明是 $\mathcal F_2$ 仍可选择 $\mathcal F_1$ 的最优解。其余功率、接入、SOC、效率和服务约束相同时，V2G 允许放电为零，故包含原只充电集合。新增工作人员接入窗口只有在允许保留旧计划、任务与成本预算一致时才满足包含关系。移动会消耗能量和时间，各地点之间一般没有包含关系。

式\eqref{evf:inclusion}是全局最优值的条件结论。当前局部搜索保留旧可行曲线，并在同一信息指标下只接受实算增益；扩展设备或接入域时，也重新验证原曲线并保留其已取得分数。这样，增加设备能力至少可维持已取得的旧计划。人员模型当前提供已声明的接线窗口、接线时长和离站时间；固定地点下，较早完成且没有额外成本的接线弱增大后续功率机会。完整车辆与人员联合路线尚未进入本批。

短时间窗可能使 V2G 无法提供放电。令 $G=e_{\rm depart}-e_{\rm initial}+E_{{\rm drive},j}>0$ 为需补入电池的能量，有 $n$ 个可接入时段，每段功率恒定，充电上限 $\bar u$。若有一个时段放电，最多剩 $n-1$ 个时段充电，因此
\begin{equation}
 \eta_c\Delta t\bar u(n-1)<G
 \quad\Longrightarrow\quad\text{该任务下任何可行曲线均不能放电。}
 \label{evf:short-window}
\end{equation}
放电还会增加所需补能，故这个条件充分。该结论针对当前按整段功率和互斥建模的时间尺度；更短控制周期应使用新的可执行窗口核查。

\subsubsection{连续曲线怎样优化，怎样检查实际辨识}

新增 \texttt{ev\_service\_milp.py} 固定终端电池能量，并加入接入、行驶、接线、离站、返程备用、效率、SOC及充放互斥。固定地点只充电时，式\eqref{si:MILP}的操作子问题是 LP；双向互斥时是 MILP。内部最小化 $-\eta$，其最小化对偶下界取负后是该仿射最大化的上界；解码可行曲线给出仿射值下界。LP 另保存对偶目标、对偶一阶条件残差和乘子符号误差；所有曲线独立重算能量、SOC及服务任务。数值界只对应当次仿射子问题。

搜索仍从恒充和操作优化生成的初值开始，以完整 AC 导数构成价格步，再重算式\eqref{si:robust-task}决定接受。几个接入点在外层比较，连续曲线没有预先枚举。两时段秩一解析例可用方向二进制精确优化平方距离；一般 AC 信息目标继续是局部多初值非凸搜索。

设计指标 $\Phi$ 反映声明工程 R/X 方向上、允许其他参数局部补偿后留下的信息。实际评价对每个候选取得完整 H+N 的白化 SSE $S_h$；给定评价真候选 $h_*$，记录
\begin{equation}
 \boxed{M=\min_{h\ne h_*}S_h-S_{h_*}.}
 \label{evf:actual-margin}
\end{equation}
$M>0$ 表示本次取得的拟合中真候选领先；共同已知协方差的 Gaussian 计量下，最大已取得对数似然的相应领先量是 $M/2$。必要拟合失败时，$M$ 记为缺失；失败保留全部候选的拒绝分数为零，不替代原拟合误差或 $M$。真候选仅供独立评价使用，选择器只读取当前历史、历史拟合、声明预测和公开设备条件。各候选线路跨 H+N 共享，逐时真实 P/Q 与同接口源同时重拟合；局部拟合结果和设计分数分别保留。

$\Phi$ 与 $M$ 的对象不同：前者平均预定方向并在局部展开点比较支路收缩，后者由实际新增数据、全部竞争拓扑和完整非线性补偿决定。强 R 可以提高平均信息，弱 X、零边结构或另一合法正边竞争仍可能未决；高代理分数而实际裕量下降的运行保留用于判断指标缺口。未来负荷来自历史均值预测，时段质量差异属于公开计量条件，最佳实际负荷时段仍需另行预测验证。

独立解析核查已通过，覆盖历史补偿、终端服务和损耗、EV/ESS连接状态等价、人员窗口正反例、旅行耗尽、纯P弱X和零梯度停留。原文设备基础来自 Al Taha、Vincent、Bitar 的 ACC 2023 第II节（无损累计能量），以及 Panda、Tindemans 的 EPSR 2024 第2节（只充电、统一接入窗口）；效率、V2G、驾驶和人员条件及本节辨识推导由本研究增加。开放全文和阅读范围保存在 \texttt{literature\_provenance.json}，两份来源分别为
\url{https://bitar.engineering.cornell.edu/papers/ACC23.pdf}、
\url{https://doi.org/10.1016/j.epsr.2024.110811}。
\subsubsection{六个独立历史的开发验证：曲线有效，优势有条件}

新开发批次先完成两个独立历史的接口与耗时检查，再按预声明条件扩至六个历史。每个历史固定相同的 H、候选和 H-only 拟合，分别检查 uniform、temporal、spatial 三种公开计量质量，共18个条件、9种方法。无EV是同长度计量诊断，未完成充电服务；其余方法采用相同终端电池任务。三个质量条件来自同一个历史，配对比较先在 H 内平均三种质量，再对6个独立 H 汇总。本批用于开发判断，未使用旧校准门槛。

独立重放核对72份保存历史拟合副本，其中24个实际唯一历史拟合；648份联合拟合副本，其中376个实际唯一拟合；162组真实新增计量响应。21份实际科学源码与冻结快照一致，全部通过。算法共42055次 AC 调用，计入历史拟合、规划和新的联合拟合，排除仿真真值生成。各条件均通过预声明接口和耗时门。

\begin{table}[htbp]
 \centering\small
 \caption{开发批次的实际 AC 裕量差。每个 H 先平均三种质量，再平均6个 H；正、负、零按 H 计。}
 \label{evf:development-pairs}
 \begin{tabular}{lrc}
  \toprule
  配对：前者减后者 & 平均 $\Delta M$ & 正/负/零 H\\
  \midrule
  地点0信息曲线--地点0恒充 & 393.008 & 5/1/0\\
  移动单向--地点0信息曲线 & 9.373 & 1/0/5\\
  地点0 V2G--地点0单向 & 497.137 & 3/3/0\\
  人员扩窗--移动单向 & $-0.128$ & 0/1/5\\
  移动单向--简单质量曲线 & 27.826 & 3/3/0\\
  \bottomrule
 \end{tabular}
\end{table}

连续曲线确由操作优化产生。地点0只充电在18条件中的8个取得超过最佳优化初值的后续价格增益；平均该项为1.390、最大5.969。地点0 V2G 在9个条件有后续增益，平均27.269、最大60.586。这些数值是 $\Phi$ 的改进，初值贡献和最终 $M$ 另存。两种 V2G 各有1个条件保留扩展前的已知可行曲线，其余17个采用信息搜索结果；随机选点在6/18条件请求不可完成服务任务的地点，保留选择失败并安全回退至地点0恒充。

移动优势在一个明确条件下出现：历史981004的 spatial 条件，地点0信息曲线取得 $M=-.3413$，移动到地点1的曲线为 $(0,6,.8421,6)$ kW，取得 $M=168.3637$。它同时补回 $.2$ kWh 驾驶耗电，达到相同末端电量。六 H 的移动增益集中在该历史，其余五 H 选择保留原地点；该结果是完整移动计划的增益，同时改变了位置、可充时段和补回驾驶的网侧电量。位置单独的作用还需匹配曲线及电量的诊断对照，稳定改善由新历史评价。

更高信息分数也会对应更差的实际裕量。弱支路历史981005的 spatial 条件，人员扩窗后选择地点2，$\Phi$ 从40.4895升至40.6236，而 $M$ 从 $-1.1466$ 降至 $-3.4514$；该条运行保留。人员扩窗使原不可完成任务的地点变为可行，本批尚未取得额外实际排错收益。人员接线时长为0，表示已准备接口的理想外部窗口；联合人员路线和实际接线耗时的收益继续另查。

本批固定初始24 kWh、末端26.85 kWh、功率上限6 kW、效率 $.95$、四个15分钟时段。固定点恒充从网侧取得3 kWh；移动点需补3.05 kWh到电池。原地点2只剩两时段，最多补2.85 kWh，所以保留操作不可行。新增窗口后的移动点仍只有三个充电时段；按式\eqref{evf:short-window}，它无法留出一个放电时段。固定点则能够生成 $(6,6,6,-5.415)$ kW 的 V2G 曲线，其网侧净电量为3.14625 kWh，转换损耗 $.29625$ kWh。设备自由度是否有用取决于服务任务和可接入时间。

开发证据支持当前方法产生可执行且有信息增益的连续曲线，也给出了固定点曲线和移动位置改善实际 AC 裕量的条件。V2G 相对单向、人员扩窗相对原窗口、信息选址相对简单质量规则的结果仍有正负差异；平均代理改善不替代实际排错和风险评价。下一步冻结新策略，使用新独立历史校准与评价；本批历史继续仅作为开发数据。
